\subsection{Macaulay 环}

定义 3.--- 称环 A 为 Macaulay 的，或称为 Macaulay 环，如果它是 Noether 的且 A-模 A 是 Macaulay 的。

例.- 1) 每个 Artin 环都是 Macaulay 环 (第 1 节, 例 1)。2) Macaulay 环没有嵌入伴随素理想 (第 2 节, 命题 2)。反之，设 A 为没有嵌入伴随素理想的维数 ≤ 1 的 Noether 环；对 A 的每个非空有限强割子集 S，A-模 A/SA 的维数 ≤ 0，从而是 Macaulay 的 (第 1 节, 例 1)；因此 A 是 Macaulay 环 (第 4 节, 定理 2)。特别地，维数 ≤ 1 的约化 Noether 环是 Macaulay 环。

\begin{enumerate}
\def\labelenumi{\arabic{enumi})}
\setcounter{enumi}{2}
\item
  维数 \(\leqslant 2\) 的正规 Noether 环是 Macaulay 环 (\(\S 1, n^{\circ} 10\), 定理 4 之前的正文)。反之，设 \(A\) 为一个 Macaulay 环，它在每个素理想处的局部环 \(A_{\mathfrak{p}}\)（素理想 \(\mathfrak{p}\) 的高度 \(\leqslant 1\)）都是整闭的；则 A 是正规的 (\(\S 1, n^{\circ} 10\), 参见定理 4)。
\item
  若 A 是 Macaulay 环，则对 A 的每个乘性子集 S，\(S^{-1}A\) 也是 Macaulay 环 (\(n^{\circ}\) 1, 注)。反之，若对 A 的每个极大理想 m，环 \(A_{m}\) 都是 Macaulay 环，则环 A 是 Macaulay 的 (\(n^{\circ}\) 1, 定义 1)。
\item
  设 A 为 Noether 环，J 为 \(A\) 的理想。为使 A/J 是 Macaulay 环，必须且只需它是 \(A\)-Macaulay 模 (第 1 节, 例 4)。
\item
  设 A 为 Noether 局部环，J 为由 A-正则序列生成的 A 的理想。为使 A/J 是 Macaulay 环，必须且只需 A 本身是 (第 3 节, 例 5 与命题 4)。
\item
  为使 Noether 局部环 A 是 Macaulay 的，必须且只需它具有一个由 A-正则序列生成的定义理想：这由第 3 节, 命题 3 以及下述事实得出：m\_A 中元素的 A-正则序列生成一个定义理想当且仅当其长度为 dim(A) (VIII, § 3, 第 2 节, 定理 1 及命题 3 的推论)。特别地，每个正则 Noether 局部环都是 Macaulay 环 (VIII, § 5, 第 2 节, 定理 1)。更一般地，正则 Noether 局部环 A 对由 A-正则序列生成的理想的商是 Macaulay 环 (例 6)。
\end{enumerate}

命题 7.-- 设 A 为 Noether 环。下列条件等价：

\begin{enumerate}
\def\labelenumi{(\roman{enumi})}
\item
  A 是 Macaulay 环；
\item
  对 Spec(A) 的每个闭子集 F，有 \(\mathrm{prof}_{\mathrm{F}}(\mathrm{A}) = \mathrm{codim}(\mathrm{F})\)；
\item
  A 的每个理想 J 都含有一个长度为 ht(J) 的 A-正则序列；
\end{enumerate}

(iii') A 的每个极大理想 m 都含有一个长度为 ht(m) 的 A-正则序列；

\begin{enumerate}
\def\labelenumi{(\roman{enumi})}
\setcounter{enumi}{3}
\tightlist
\item
  对 A 的每个理想 J，有 \(\mathrm{Ext}_A^i (\mathrm{A} / \mathrm{J},\mathrm{A}) = 0\) 对 \(i <   \mathrm{ht}(\mathrm{J})\) 成立
\end{enumerate}

(iv') 对 A 的每个极大理想 m，有 Ext \(_{A}^{i}\) (A/m, A) = 0 对 i \textless{} ht(m) 成立；

\begin{enumerate}
\def\labelenumi{(\alph{enumi})}
\setcounter{enumi}{21}
\tightlist
\item
  对 A 的每个素理想 p 以及 \(A_{p}\) 的每个由 \(A_{p}\) 的极大割序列生成的理想 J，有 \(e_{\mathrm{J}}(\mathrm{A}_{\mathfrak{p}})=\mathrm{long}(\mathrm{A}_{\mathfrak{p}}/\mathrm{JA}_{\mathfrak{p}})\)；
\end{enumerate}

(v') 对 A 的每个极大理想 m，存在 \(A_{m}\) 的一个理想 J，它由 \(A_{m}\) 的一个极大割序列生成，且满足 \(e_{\mathrm{J}}(A_{\mathrm{m}})=\mathrm{long}(\mathrm{A}_{\mathrm{m}}/\mathrm{JA}_{\mathrm{m}})\)；

\begin{enumerate}
\def\labelenumi{(\roman{enumi})}
\setcounter{enumi}{5}
\tightlist
\item
  (Macaulay-Cohen 判别法) 对 A 的每个使得由 S 生成的理想 G 具有高度 Card(S) 的有限子集 S，A-模 A/G 都没有嵌入伴随素理想。
\end{enumerate}

(i) 与 (ii) 的等价性由第 1 节, 命题 1 的推论得出。由 § 1, 第 4 节, 定理 2 以及深度的定义，条件 (iii) 与 (iv)（相应地 (iii') 与 (iv')）意味着对 A 的每个理想（相应地每个极大理想）J 有 \(\mathrm{prof}_{\mathrm{A}}(\mathrm{J};\mathrm{A})\geqslant \mathrm{ht}(\mathrm{J})\)。于是有

\[
(i) \Leftrightarrow (ii) \Rightarrow (iii) \Leftrightarrow (iv) \Rightarrow (iii') \Leftrightarrow (iv').
\]

但 (iv') 蕴含：对 A 的每个极大理想 m，\(\operatorname{Ext}_{A_{m}}^{i}(\kappa(m), A_{m}) = 0\) 对 \(i < \dim(A_{m})\) 成立，于是 \(\operatorname{prof}(A_{m}) \geqslant \dim(A_{m})\)，从而 A 是 Macaulay 环。

(i)、(v) 与 (v') 的等价性由第 3 节, 定理 1 得出，而 (i) 与 (vi) 的等价性由第 4 节, 定理 2 得出。

